\subsection{Elementary Divisors and Invariant Factors}

\begin{definition} \label{def:elementary-divisors-pid-module} \hlt{Elementary Divisors}\\
Set $\operatorname{ElemDiv}(0)\defeq\emptyset$. Let $0\neq M$ be f.g. torsion. Choose its intrinsic primary components by Theorem \ref{thm:primary-decomposition-pid-modules}.\\
Choose their intrinsic exponent chains by Theorem \ref{thm:cyclic-decomposition-primary-pid}. Suppose the resulting cyclic summands have annihilators $\langle p_i^{e_{i,j}}\rangle$.\\
Define the \hlt{elementary-divisor multiset} by $\operatorname{ElemDiv}(M)\defeq\{p_i^{e_{i,j}}\}_{i,j}$, with entries taken up to associates.
\end{definition}

\begin{theorem} \label{thm:primary-cyclic-decomposition-pid} \hlt{Primary Cyclic Decomposition}\\
Let $M$ be a f.g. torsion $R$-module.
\begin{enumerate}[label=(\roman*)]
\setlength{\itemsep}{0pt}
\item There is a finite set $I$, together with distinct nonassociate primes $(p_i)_{i\in I}$ and $(k_i)_{i\in I}\in\N^I$.\\
There are exponents $e_{i,1}\geq\cdots\geq e_{i,k_i}\geq1$ and vectors $v_{i,j}\in M$ $(i\in I,\ 1\leq j\leq k_i)$. These data satisfy $M=\bigoplus\limits_{i\in I}\bigoplus\limits_{j=1}^{k_i}\langle\!\langle v_{i,j}\rangle\!\rangle_R$.\\
Their annihilators satisfy $\operatorname{Ann}_R(v_{i,j})=\langle p_i^{e_{i,j}}\rangle$ $(i\in I,\ 1\leq j\leq k_i)$. For $M=\{0_M\}$, take $I=\emptyset$.
\item If $N$ is f.g. torsion, then $M\cong N\Longleftrightarrow\operatorname{ElemDiv}(M)=\operatorname{ElemDiv}(N)$.
\end{enumerate}
\begin{proof}
\begin{enumerate}[label=(\roman*)]
\setlength{\itemsep}{0pt}
\item[$(i)$] If $M=\{0_M\}$, the empty direct sum has the asserted form. Suppose $M\neq\{0_M\}$.\\
By Theorem \ref{thm:primary-decomposition-pid-modules}, there is a finite family of f.g. $p_i$-primary components. It satisfies $M=\bigoplus\limits_{i\in I}M_{p_i}$.\\
By Theorem \ref{thm:cyclic-decomposition-primary-pid}(i), each $M_{p_i}$ has an internal cyclic decomposition. Its exponent chain and annihilators have the asserted form.\\
Substitution into the primary decomposition proves the displayed internal direct sum.
\item[$(ii)$] ($\Rightarrow$) Suppose $\Phi:M\xrightarrow{\sim}N$. The map $\Phi$ identifies the intrinsic primary components by Theorem \ref{thm:primary-decomposition-pid-modules}(iii).\\
Their exponent multisets agree by Theorem \ref{thm:cyclic-decomposition-primary-pid}(ii), so $\operatorname{ElemDiv}(M)=\operatorname{ElemDiv}(N)$.
\item[$(ii)$] ($\Leftarrow$) Suppose $\operatorname{ElemDiv}(M)=\operatorname{ElemDiv}(N)$. Set $C\defeq\bigoplus\limits_{q\in\operatorname{ElemDiv}(M)}R/\langle q\rangle$. Then $M\cong C\cong N$ by part (i) and Theorem \ref{thm:cyclic-modules-over-pid}(i).
\end{enumerate}
\noindent
\hlt{Structural Commentary}: Intrinsic primary kernels and exponent layers make the prime-power blocks independent of choices.\\
\hlt{Key Insight}: The elementary-divisor multiset determines the common quotient-sum model.
\end{proof}
\end{theorem}

\begin{proposition} \label{prop:elementary-divisors-direct-sums} \hlt{Elementary Divisors of a Direct Sum}\\
Let $A$ and $B$ be f.g. torsion, and let $\uplus$ denote multiset union. Then $\operatorname{ElemDiv}(A\oplus B)=\operatorname{ElemDiv}(A)\uplus\operatorname{ElemDiv}(B)$.
\begin{proof}
Choose primary cyclic decompositions of $A$ and $B$ by Theorem \ref{thm:primary-cyclic-decomposition-pid}(i). Write $A=\bigoplus\limits_{i\in I}A_i$ and $B=\bigoplus\limits_{j\in J}B_j$.\\
Then $A\oplus B=(\bigoplus\limits_{i\in I}A_i)\oplus(\bigoplus\limits_{j\in J}B_j)$. This is a primary cyclic decomposition, including when $I$ or $J$ is empty.\\
By Definition \ref{def:elementary-divisors-pid-module}, its prime-power annihilators form the required multiset. That multiset is $\operatorname{ElemDiv}(A)\uplus\operatorname{ElemDiv}(B)$.\\
\hlt{Structural Commentary}: Concatenating primary cyclic decompositions concatenates their annihilator multisets.\\
\hlt{Key Insight}: Multiset union retains repeated prime-power blocks.
\end{proof}
\end{proposition}

\begin{definition} \label{def:invariant-factors-pid-module} \hlt{Invariant Factors}\\
Let $M$ be f.g. torsion. If $\operatorname{ElemDiv}(M)=\emptyset$, define $\operatorname{Inv}(M)\defeq()$. Suppose $\operatorname{ElemDiv}(M)\neq\emptyset$.\\
Group it as $\{p_i^{e_{i,j}}:i\in I,\ 1\leq j\leq k_i\}$; $p_i$ are distinct nonassociate primes. Order each exponent row by $e_{i,1}\geq\cdots\geq e_{i,k_i}\geq1$.\\
Set $m\defeq\max\limits_{i\in I}k_i$ and extend every row by $e_{i,j}\defeq0$ for $k_i<j\leq m$. Set $d_j\defeq\prod\limits_{i\in I}p_i^{e_{i,j}}$ $(1\leq j\leq m)$.\\
\hlt{Invariant-factor tuple} is $\operatorname{Inv}(M)\defeq(d_1,\ldots,d_m)$, with entries taken up to associates. Entries are nonunits and satisfy $d_m\mid\cdots\mid d_1$.
\end{definition}

\begin{method} \label{meth:elementary-divisors-to-invariant-factors} \hlt{From Elementary Divisors to Invariant Factors}
\begin{enumerate}[label=(\roman*)]
\setlength{\itemsep}{0pt}
\item If $\operatorname{ElemDiv}(M)=\emptyset$, return $\operatorname{Inv}(M)=()$. Otherwise, group the input multiset by associate prime.
\item Write the $i$th group as $p_i^{e_{i,1}},\ldots,p_i^{e_{i,k_i}}$ with $e_{i,1}\geq\cdots\geq e_{i,k_i}$. Pad every exponent row by zeros to length $m\defeq\max\limits_i k_i$.
\item Form $d_j\defeq\prod\limits_i p_i^{e_{i,j}}$ for $1\leq j\leq m$.
\item Return $\operatorname{Inv}(M)=(d_1,\ldots,d_m)$ and verify $d_{j+1}\mid d_j$ from $e_{i,j+1}\leq e_{i,j}$.\\
Recover the input by extracting its maximal prime-power factors from each $d_j$.
\end{enumerate}
\end{method}

\begin{theorem} \label{thm:pid-invariant-factor-decomposition} \hlt{Invariant Factor Decomposition}\\
Let $M$ be f.g. torsion and write $\operatorname{Inv}(M)=(d_1,\ldots,d_m)$, where $m\in\N_0$.
\begin{enumerate}[label=(\roman*)]
\setlength{\itemsep}{0pt}
\item $M\cong\bigoplus\limits_{j=1}^mR/\langle d_j\rangle$. If $m\geq1$, then $d_j\notin R^\times$ $(1\leq j\leq m)$ and $d_{j+1}\mid d_j$ $(1\leq j<m)$.
\item If $N$ is f.g. torsion, then $M\cong N\Longleftrightarrow\operatorname{Inv}(M)=\operatorname{Inv}(N)$.
\end{enumerate}
\begin{proof}
\begin{enumerate}[label=(\roman*)]
\setlength{\itemsep}{0pt}
\item[$(i)$] If $m=0$, then $\operatorname{ElemDiv}(M)=\emptyset$ by Definition \ref{def:invariant-factors-pid-module}.\\
Hence $M=\{0_M\}$ by Theorem \ref{thm:primary-cyclic-decomposition-pid}(i), so the empty direct sum proves the claim. Suppose $m\geq1$.\\
By Definition \ref{def:invariant-factors-pid-module}, $\operatorname{ElemDiv}(M)=\{p_i^{e_{i,j}}:i\in I,\ 1\leq j\leq k_i\}$. Pad exponent rows by zeros, set $d_j=\prod\limits_{i\in I}p_i^{e_{i,j}}$.\\
By Theorem \ref{thm:primary-cyclic-decomposition-pid}(i), $M$ is a direct sum of the associated primary cyclic blocks.\\
Each block is isomorphic to its quotient model by Theorem \ref{thm:cyclic-modules-over-pid}(i).\\
Hence $M\cong\bigoplus\limits_{i\in I}\bigoplus\limits_{j=1}^{k_i}R/\langle p_i^{e_{i,j}}\rangle$. For $1\leq j\leq m$, set $C_j\defeq\bigoplus\limits_{\substack{i\in I\\e_{i,j}>0}}R/\langle p_i^{e_{i,j}}\rangle$. By finite reindexing, $M\cong\bigoplus\limits_{j=1}^mC_j$.\\
The nonunit orders in $C_j$ are pairwise coprime. By Theorem \ref{thm:cyclic-composition-coprime-orders}(i)--(ii), $C_j$ is cyclic and has order associated to $d_j$.\\
Hence $C_j\cong R/\langle d_j\rangle$ by Theorem \ref{thm:cyclic-modules-over-pid}(i). Substitution proves the asserted decomposition.\\
The nonunit and divisibility properties hold by Definition \ref{def:invariant-factors-pid-module}.
\item[$(ii)$] ($\Rightarrow$) Suppose $M\cong N$. Then $\operatorname{ElemDiv}(M)=\operatorname{ElemDiv}(N)$ by Theorem \ref{thm:primary-cyclic-decomposition-pid}(ii); $\operatorname{Inv}(M)=\operatorname{Inv}(N)$ by Definition \ref{def:invariant-factors-pid-module}.
\item[$(ii)$] ($\Leftarrow$) Suppose $\operatorname{Inv}(M)=\operatorname{Inv}(N)$. Factoring  invariant factors recovers equal prime-power multisets by Theorem \ref{thm:pid-is-ufd}(ii).\\
Thus $\operatorname{ElemDiv}(M)=\operatorname{ElemDiv}(N)$. Hence $M\cong N$ by Theorem \ref{thm:primary-cyclic-decomposition-pid}(ii).
\end{enumerate}
\noindent
\hlt{Structural Commentary}: Coprime cyclic composition packs a single block from each prime into a divisibility column.\\
\hlt{Key Insight}: Factoring the invariant factors reverses the column packing.
\end{proof}
\end{theorem}

\begin{example} \label{ex:elementary-divisors-to-invariant-factors-polynomial-module} \hlt{Packing Elementary Divisors into Invariant Factors}\\
Let $R=\mathbb{R}[z]$ and $\operatorname{ElemDiv}(M)=\{(z-1)^3,z-1,(z^2+1)^2,z^2+1\}$. The polynomials $z-1$ and $z^2+1$ are nonassociate primes in $R$.\\
By Method \ref{meth:elementary-divisors-to-invariant-factors}, $d_1=(z-1)^3(z^2+1)^2$ and $d_2=(z-1)(z^2+1)$. Thus $\operatorname{Inv}(M)=(d_1,d_2)$ with $d_2\mid d_1$.\\
By Theorem \ref{thm:pid-invariant-factor-decomposition}(i), $M\cong R/\langle d_1\rangle\oplus R/\langle d_2\rangle$.\\
\hlt{Structural Commentary}: Separate prime rows become the two invariant-factor columns.\\
\hlt{Key Insight}: Align exponents by height before multiplying across primes.
\end{example}

\begin{example} \label{ex:elementary-divisors-recombine-to-invariant-factors} \hlt{Recombining Elementary Divisors}\\
As a $\Z$-module, let $G\defeq\Z/12\Z\oplus\Z/18\Z$. By the primary decompositions of its cyclic summands, $\operatorname{ElemDiv}(G)=\{2^2,2,3^2,3\}$.\\
The exponent rows are $(2,1)$ for both primes, so $d_1=2^2\cdot3^2=36$ and $d_2=2\cdot3=6$.\\
Thus $\operatorname{Inv}(G)=(36,6)$ and $G\cong\Z/36\Z\oplus\Z/6\Z$ by Theorem \ref{thm:pid-invariant-factor-decomposition}(i), with $6\mid36$.\\
\hlt{Structural Commentary}: The $2$- and $3$-primary rows recombine into the divisibility chain $6\mid36$.\\
\hlt{Key Insight}: Matching exponent ranks determines which prime powers multiply together.
\end{example}

\begin{corollary} \label{cor:general-finitely-generated-pid-module-classification} \hlt{Classification of Finitely Generated PID Modules}\\
Let $M$ and $N$ be f.g. $R$-modules. Set $T_X\defeq X_{\operatorname{tor}}$ and $\rho_X\defeq\operatorname{rk}_R(X/T_X)$ $(X\in\{M,N\})$.
\begin{enumerate}[label=(\roman*)]
\setlength{\itemsep}{0pt}
\item There are $r,m\in\N_0$ and nonunits $d_1,\ldots,d_m\in R$ with $d_{j+1}\mid d_j$ $(1\leq j<m)$. For these data, $M\cong R^r\oplus\bigoplus\limits_{j=1}^mR/\langle d_j\rangle$.
\item $M\cong N\Longleftrightarrow[\rho_M=\rho_N\land\operatorname{Inv}(T_M)=\operatorname{Inv}(T_N)]$.
\item $M\cong N\Longleftrightarrow[\rho_M=\rho_N\land\operatorname{ElemDiv}(T_M)=\operatorname{ElemDiv}(T_N)]$.
\end{enumerate}
\begin{proof}
\begin{enumerate}[label=(\roman*)]
\setlength{\itemsep}{0pt}
\item[$(i)$] By Theorem \ref{thm:finitely-generated-pid-free-plus-torsion}(i), choose f.g. free $F\leq M$ with $M=T_M\oplus F$, $F\cong M/T_M$, $r\defeq\operatorname{rk}_R(F)\in\N_0$; $F\cong R^r$ by Proposition \ref{prop:free-module-coordinate-model-and-isomorphism}(i).\\
Coordinate projection $M\twoheadrightarrow T_M$ maps finite generating set of $M$ onto generating set of $T_M$, so $T_M$ is f.g.\\
Write $\operatorname{Inv}(T_M)=(d_1,\ldots,d_m)$ with $m\in\N_0$. By Theorem \ref{thm:pid-invariant-factor-decomposition}(i), these data have the asserted properties for $T_M$.\\
Substitution proves the claim.
\item[$(ii)$] ($\Rightarrow$) Suppose $\Phi:M\xrightarrow{\sim}N$ and fix $x\in T_M$. Choose $a\in R\setminus\{0_R\}$ with $ax=0_M$. Then $a\Phi(x)=0_N$, so $\Phi(T_M)\subseteq T_N$.\\
The same argument for $\Phi^{-1}$ proves $\Phi(T_M)=T_N$. Thus $\Phi$ restricts to $T_M\xrightarrow{\sim}T_N$ and induces $M/T_M\xrightarrow{\sim}N/T_N$.\\
Hence $\rho_M=\rho_N$ by Proposition \ref{prop:free-module-coordinate-model-and-isomorphism}(ii). The invariant-factor tuples agree by Theorem \ref{thm:pid-invariant-factor-decomposition}(ii).
\item[$(ii)$] ($\Leftarrow$) Suppose $\rho_M=\rho_N$ and $\operatorname{Inv}(T_M)=\operatorname{Inv}(T_N)$. Then $M/T_M\cong N/T_N$ by Proposition \ref{prop:free-module-coordinate-model-and-isomorphism}(ii).\\
Torsion modules satisfy $T_M\cong T_N$ by Theorem \ref{thm:pid-invariant-factor-decomposition}(ii). As $M\cong T_M\oplus M/T_M$, $N\cong T_N\oplus N/T_N$ by Theorem \ref{thm:finitely-generated-pid-free-plus-torsion}(i), $M\cong N$.
\item[$(iii)$] ($\Rightarrow$) Suppose $M\cong N$. Then $\rho_M=\rho_N$ by part (ii). The given module isomorphism induces $T_M\cong T_N$.\\
Hence $\operatorname{ElemDiv}(T_M)=\operatorname{ElemDiv}(T_N)$ by Theorem \ref{thm:primary-cyclic-decomposition-pid}(ii).
\item[$(iii)$] ($\Leftarrow$) Suppose $\rho_M=\rho_N$ and $\operatorname{ElemDiv}(T_M)=\operatorname{ElemDiv}(T_N)$. Then $M/T_M\cong N/T_N$ by Proposition \ref{prop:free-module-coordinate-model-and-isomorphism}(ii).\\
The torsion modules satisfy $T_M\cong T_N$ by Theorem \ref{thm:primary-cyclic-decomposition-pid}(ii).\\
The free-plus-torsion decompositions therefore imply $M\cong N$ by Theorem \ref{thm:finitely-generated-pid-free-plus-torsion}(i).
\end{enumerate}
\noindent
\hlt{Structural Commentary}: The intrinsic torsion summand carries invariant factors; the free quotient carries the rank.\\
\hlt{Key Insight}: Isomorphism requires equality of both independent data sets.
\end{proof}
\end{corollary}

\begin{method} \label{meth:smith-form-module-presentation} \hlt{Reading a Module from Smith Form}\\
Let $R^n\xrightarrow{A}R^m\to M\to0$ be a finite presentation. Set $D\defeq\operatorname{diag}(d_1,\ldots,d_r,0,\ldots,0)$.\\
Suppose $A\sim_{\mathrm{eq}}D$ and $0\leq r\leq\min\{m,n\}$. If $r\geq1$, suppose $(d_1,\ldots,d_r)\in(R\setminus\{0_R\})^r$ and $d_1\mid\cdots\mid d_r$.
\begin{enumerate}[label=(\roman*)]
\setlength{\itemsep}{0pt}
\item Record the torsion-coordinate quotients $(R/\langle d_i\rangle)_{i=1}^r$ and the free-coordinate module $R^{m-r}$.
\item Delete the zero quotients indexed by $d_i\in R^\times$. Remaining list, together with $R^{m-r}$, is decomposition from Corollary \ref{cor:smith-normal-form-presents-pid-modules}.
\item If every $d_i$ is a unit, return $\operatorname{Inv}(M_{\operatorname{tor}})=()$. Otherwise, if nonunit entries are $d_s,\ldots,d_r$, return $\operatorname{Inv}(M_{\operatorname{tor}})=(d_r,\ldots,d_s)$.
\item In nonempty case, factor each $d_i$ into pairwise nonassociate maximal prime powers. Retain multiplicities for $\operatorname{ElemDiv}(M_{\operatorname{tor}})$.
\end{enumerate}
\hlt{Key Insight}: Reverse the nonunit Smith entries for the repository's largest-first invariant-factor tuple.
\end{method}

\begin{perspective} \label{per:invariant-factors-order-and-characteristic-polynomial} \hlt{Largest Invariant Factor and Product}\\
Let $0\neq M$ be f.g. torsion with $\operatorname{Inv}(M)=(d_1,\ldots,d_m)$. Decomposition $M\cong\bigoplus\limits_{j=1}^mR/\langle d_j\rangle$ holds by Theorem \ref{thm:pid-invariant-factor-decomposition}(i).\\
Thus $\operatorname{Ann}_R(M)=\langle\operatorname{lcm}(d_1,\ldots,d_m)\rangle$ by Proposition \ref{prop:orders-direct-sum-lcm}. Since $d_m\mid\cdots\mid d_1$, this ideal is $\langle d_1\rangle$.\\
The column construction in Definition \ref{def:invariant-factors-pid-module} uses each elementary divisor once. Hence $\prod\limits_{j=1}^md_j\sim\prod\limits_{q\in\operatorname{ElemDiv}(M)}q$.\\
Thus the largest invariant factor controls the annihilator, whereas the total product records every elementary-divisor block.
\end{perspective}

\begin{theorem} \label{thm:cyclic-torsion-module-characterization} \hlt{Characterizing Cyclic Torsion Modules}\\
Let $n\in\N$, let $0\neq M$ be f.g. torsion, and let $p_1,\ldots,p_n$ be distinct nonassociate primes.\\
Suppose $e_1,\ldots,e_n\in\N$ and $o(M)\sim\mu\defeq\prod\limits_{i=1}^np_i^{e_i}$. Set $I\defeq\{1,\ldots,n\}$. The following are equivalent:
\begin{enumerate}[label=(\roman*)]
\setlength{\itemsep}{0pt}
\item $M$ is cyclic.
\item The primary components $M_{p_i}$ are cyclic $(i\in I)$.
\item $\operatorname{ElemDiv}(M)=\{p_i^{e_i}:i\in I\}$.
\item $M\cong R/\langle\mu\rangle$.
\end{enumerate}
\begin{proof}
\begin{enumerate}[label=(\roman*)]
\setlength{\itemsep}{0pt}
\item[$(i)$]$\Rightarrow (ii)$ Suppose $M=\langle\!\langle x\rangle\!\rangle_R$ and choose the order representative $o(x)=\mu$. By Theorem \ref{thm:cyclic-composition-coprime-orders}(iii), choose $x_1,\ldots,x_n\in M$.\\
These vectors satisfy $M=\bigoplus\limits_{i=1}^n\langle\!\langle x_i\rangle\!\rangle_R$. These vectors satisfy $o(x_i)\sim p_i^{e_i}$ $(i\in I)$.\\
Hence $M_{p_i}=\langle\!\langle x_i\rangle\!\rangle_R$ by Theorem \ref{thm:primary-decomposition-pid-modules}(iii), so every primary component is cyclic.\\
\item[$(ii)$]$\Rightarrow (iii)$ By Theorem \ref{thm:primary-decomposition-pid-modules}(i)--(ii), the primary decomposition is indexed by $I$. It satisfies $M=\bigoplus\limits_{i=1}^nM_{p_i}$ and $o(M_{p_i})\sim p_i^{e_i}$ $(i\in I)$.\\
Choose generators $(y_i)_{i\in I}$ with $M_{p_i}=\langle\!\langle y_i\rangle\!\rangle_R$. Then $\operatorname{Ann}_R(y_i)=\operatorname{Ann}_R(M_{p_i})=\langle p_i^{e_i}\rangle$ $(i\in I)$.\\
Hence $M_{p_i}\cong R/\langle p_i^{e_i}\rangle$ by Theorem \ref{thm:cyclic-modules-over-pid}(i). Thus $\operatorname{ElemDiv}(M_{p_i})=\{p_i^{e_i}\}$ by Definition \ref{def:elementary-divisors-pid-module}.\\
By Proposition \ref{prop:elementary-divisors-direct-sums}, iterated over the primary components, $\operatorname{ElemDiv}(M)=\{p_i^{e_i}:i\in I\}$.\\
\item[$(iii)$]$\Rightarrow (iv)$ By Definition \ref{def:invariant-factors-pid-module}, $\operatorname{Inv}(M)$ has unique entry $d_1$. It satisfies $d_1\sim\prod\limits_{i=1}^np_i^{e_i}\sim\mu$; $M\cong R/\langle d_1\rangle=R/\langle\mu\rangle$ by Theorem \ref{thm:pid-invariant-factor-decomposition}(i).\\
\item[$(iv)$]$\Rightarrow (i)$ The module $R/\langle\mu\rangle$ is generated by $1+\langle\mu\rangle$. Hence its image under any isomorphism to $M$ generates $M$.
\end{enumerate}
\hlt{Structural Commentary}: Coprime primary generators combine into single generator exactly when each prime contributes a single block.\\
\hlt{Key Insight}: Repeated same-prime blocks obstruct cyclicity.
\end{proof}
\end{theorem}

\begin{definition} \label{def:indecomposable-module} \hlt{Indecomposable Module}\\
A nonzero $R$-module $M$ is \hlt{indecomposable} if $(\forall\,A,B\leq M)\,[M=A\oplus B\Rightarrow A=\{0_M\}\lor B=\{0_M\}]$.
\end{definition}

\begin{theorem} \label{thm:indecomposable-fg-torsion-pid-modules} \hlt{Indecomposable Finitely Generated Torsion Modules}\\
Let $0\neq M$ be f.g. torsion. The following are equivalent:
\begin{enumerate}[label=(\roman*)]
\setlength{\itemsep}{0pt}
\item $M$ is indecomposable.
\item $M$ is primary cyclic.
\item $(\exists\,p\in R)(\exists\,e\in\N)$ such that $p$ is prime and $\operatorname{ElemDiv}(M)=\{p^e\}$.
\end{enumerate}
\begin{proof}
\begin{enumerate}[label=(\roman*)]
\setlength{\itemsep}{0pt}
\item[$(i)$]$\Rightarrow (ii)$ By Theorem \ref{thm:primary-cyclic-decomposition-pid}(i), $M$ has a finite nonempty primary cyclic decomposition. Its summands are nonzero.\\
After finite reindexing, write $M\cong\bigoplus\limits_{\ell=1}^mC_\ell$ with $m\in\N$. Transport the coordinate summands into $M$.\\
If $m\geq2$, a coordinate summand and the sum of the others form a nontrivial decomposition, contrary to indecomposability.\\
Hence $m=1$, so $M$ is primary cyclic.
\item[$(ii)$]$\Rightarrow (iii)$ Choose a prime $p\in R$ and $e\in\N$ such that $M$ is cyclic, $p$-primary, and $o(M)\sim p^e$.\\
Then $M\cong R/\langle p^e\rangle$ by Theorem \ref{thm:cyclic-modules-over-pid}(i). Thus $\operatorname{ElemDiv}(M)=\{p^e\}$ by Definition \ref{def:elementary-divisors-pid-module}.
\item[$(iii)$]$\Rightarrow (i)$ Suppose $M=A\oplus B$. Coordinate projections map a finite generating set of $M$ onto generating sets of $A$ and $B$.\\
Hence both summands are f.g.; they are torsion as submodules of $M$.\\
By Proposition \ref{prop:elementary-divisors-direct-sums}, $\{p^e\}=\operatorname{ElemDiv}(A)\uplus\operatorname{ElemDiv}(B)$. Thus $\operatorname{ElemDiv}(A)=\emptyset$ or $\operatorname{ElemDiv}(B)=\emptyset$.\\
The corresponding summand is zero by Theorem \ref{thm:primary-cyclic-decomposition-pid}(i), so $M$ is indecomposable.
\end{enumerate}
\hlt{Structural Commentary}: Transported primary cyclic summands split $M$ unless the elementary-divisor multiset is a singleton.\\
\hlt{Key Insight}: A singleton elementary divisor is exactly an indecomposable block.
\end{proof}
\end{theorem}
